% Introduccion: sin numero de capitulo, hoja nueva.
% Fuente: plan aprobado, secciones 1 (Introduccion) y 2 (Planteamiento y
% justificacion del problema).
\seccionsinnumero{Introducción}

Una imagen hiperespectral, conocida por sus siglas en inglés HSI, es una representación tridimensional de una escena, conocida como cubo de datos, en la que cada píxel no almacena únicamente tres valores de color, rojo, verde y azul, como una imagen convencional, sino cientos de valores de reflectancia medidos en bandas espectrales estrechas y contiguas del espectro electromagnético. De esta manera, cada píxel contiene una \textit{firma espectral} completa, es decir, una curva que describe cómo el material presente en ese punto refleja o absorbe la energía a lo largo de un amplio rango de longitudes de onda, que típicamente abarca desde el visible hasta el infrarrojo de onda corta \parencite{bhargava2024hyperspectral}.

Desde el punto de vista de la adquisición, estas imágenes se capturan mediante sensores hiperespectrales embarcados en plataformas satelitales, aéreas o de laboratorio, que descomponen la luz incidente en sus componentes espectrales a través de elementos dispersores como prismas o redes de difracción. Los esquemas más comunes de captura son el barrido lineal, denominado \textit{pushbroom} y \textit{whiskbroom}, en el que el sensor registra franjas de la escena de forma secuencial a medida que la plataforma avanza, y los sistemas de plano focal que capturan bandas espectrales sucesivas de una misma escena. El resultado es el cubo hiperespectral, con dos dimensiones espaciales y una dimensión espectral \parencite{bhargava2024hyperspectral}.

La principal ventaja de una imagen hiperespectral frente a una imagen RGB o multiespectral convencional radica en su resolución espectral. Mientras que una imagen RGB resume toda la información en tres bandas anchas, insuficientes para distinguir materiales con colores aparentes similares, la imagen hiperespectral proporciona una firma espectral detallada que actúa como una huella dactilar del material, permitiendo diferenciar coberturas y sustancias que resultan indistinguibles al ojo humano o a una cámara ordinaria. Esta riqueza espectral habilita la identificación y clasificación de materiales con un nivel de detalle imposible de alcanzar con la información puramente cromática \parencite{bhargava2024hyperspectral}. Gracias a ello, la imagen hiperespectral se aplica hoy en agricultura de precisión, monitoreo ambiental, geología y minería, seguridad y defensa, control de calidad de alimentos y diagnóstico médico, entre muchos otros campos \parencite{bhargava2024hyperspectral}.

Dentro de la teledetección, una de las tareas de mayor interés es la \textit{detección de cambios}, conocida en inglés como \textit{change detection}, que consiste en comparar dos o más imágenes de una misma escena adquiridas en instantes de tiempo distintos, en un análisis bitemporal o multitemporal, con el fin de identificar qué zonas se han modificado. Sus aplicaciones incluyen el seguimiento de la deforestación, la expansión urbana, los desastres naturales y los cambios en el uso del suelo. En este contexto, la Detección de Cambios de Vocabulario Abierto, denominada en inglés \textit{Open Vocabulary Change Detection} u OVCD, busca no solo localizar dónde ocurrió un cambio, sino también describirlo y etiquetarlo con lenguaje natural libre, sin restringirse a un conjunto cerrado de clases definidas de antemano ni requerir un entrenamiento supervisado específico para cada tipo de cambio.

El propósito del enfoque de vocabulario abierto es superar la principal limitación de los métodos tradicionales de detección de cambios, los cuales quedan restringidos a las categorías presentes en sus datos de entrenamiento, lo que compromete su capacidad de generalización a escenarios reales \parencite{zhu2025semanticcd}. El estado del arte reciente en OVCD sobre teledetección se apoya en modelos de visión y lenguaje a gran escala para dotar de conocimiento semántico a la tarea: propuestas como \textit{Semantic-CD} incorporan el conocimiento implícito de modelos preentrenados de visión-lenguaje para mejorar la generalización entre categorías en la detección semántica de cambios \parencite{zhu2025semanticcd}. Esta línea se enmarca en el auge más amplio de la segmentación semántica de vocabulario abierto, impulsada por modelos como CLIP y su capacidad de clasificación \textit{zero-shot} \parencite{radford2021clip}.

El presente trabajo busca articular una cadena de modelos fundacionales, cada uno preentrenado a gran escala y reutilizado sin necesidad de reentrenar un clasificador de cambios que puede requerir cómputo dedicado. En primer lugar, HyperSIGMA, un modelo fundacional de comprensión hiperespectral basado en \textit{Vision Transformers}, se empleará como extractor de características para obtener representaciones profundas de cada imagen del par bitemporal \parencite{wang2024hypersigma}. A partir de estas representaciones se calculará un mapa de cambio mediante similitud coseno. Posteriormente, el modelo Segment Anything, conocido como SAM, un modelo fundacional de segmentación \textit{promptable} capaz de transferirse a nuevas distribuciones de imágenes de forma \textit{zero-shot}, refinará las regiones de cambio en máscaras precisas a nivel de píxel \parencite{kirillov2023sam}. Finalmente, CLIP, sigla de \textit{Contrastive Language-Image Pre-training}, aportará la capa de vocabulario abierto: al haber aprendido un espacio de representación compartido entre imagen y texto mediante aprendizaje contrastivo, permitirá etiquetar cada región de cambio contra un vocabulario textual libre por similitud semántica, sin clases predefinidas \parencite{radford2021clip}. De forma complementaria, la generación de descripciones en lenguaje natural de las regiones se buscará abordar con modelos de visión-lenguaje de tipo \textit{captioning} como BLIP \parencite{li2022blip}, aprovechando conjuntos de datos específicos para el dominio hiperespectral como HyperCap \parencite{das2025hypercap}.

El eje metodológico común a esta cadena es el paradigma \textit{zero-shot}: la capacidad de los modelos fundacionales de resolver tareas sobre categorías no vistas durante su entrenamiento, transfiriendo el conocimiento adquirido en grandes corpus a un dominio nuevo sin ajuste supervisado adicional. Este paradigma es precisamente lo que hace viable la detección de cambios de vocabulario abierto. No obstante, si bien la detección de cambios de vocabulario abierto se ha desarrollado sobre imágenes ópticas RGB y multiespectrales de teledetección \parencite{zhu2025semanticcd}, y la clasificación \textit{zero-shot} de imágenes hiperespectrales con CLIP ha comenzado a explorarse recientemente \parencite{pang2025special}, hasta la fecha no se ha abordado la detección de cambios de vocabulario abierto sobre imágenes hiperespectrales bitemporales, lo que constituye la brecha que motiva este trabajo.

Con base en lo anterior, esta investigación se centra en el diseño e implementación de una cadena de procesamiento, o \textit{pipeline}, \textit{zero-shot} que, a partir de un par de imágenes hiperespectrales de una misma escena tomadas en tiempos distintos, detecte las regiones de cambio y las describa con vocabulario libre, integrando modelos fundacionales de comprensión hiperespectral, segmentación y visión-lenguaje. El propósito es aprovechar la riqueza espectral de las imágenes hiperespectrales, hasta ahora poco explotada en este tipo de tareas, para lograr una caracterización semántica de los cambios más rica y generalizable que la que ofrecen los métodos supervisados de clases cerradas.

% --- Plan, seccion 2: Planteamiento y justificacion del problema ---

En los últimos años, los modelos fundacionales de visión por computador han transformado las tareas de segmentación y de comprensión semántica de imágenes. Los modelos de segmentación guiados por indicaciones, o \textit{prompts}, como Segment Anything, conocido como SAM, permiten obtener máscaras precisas de objetos sin haber sido entrenados para categorías específicas \parencite{kirillov2023sam}, mientras que los modelos de visión y lenguaje como CLIP habilitan la clasificación de vocabulario abierto al alinear imágenes y texto en un mismo espacio semántico \parencite{radford2021clip}. Sobre esta base, la detección semántica de cambios con vocabulario abierto, denominada en inglés \textit{Open-Vocabulary Semantic Change Detection}, ha comenzado a consolidarse en teledetección: métodos recientes como \textit{Semantic-CD} identifican y describen los cambios entre imágenes de una misma escena sin limitarse a un conjunto cerrado de clases definidas de antemano \parencite{zhu2025semanticcd}. No obstante, todos estos desarrollos operan sobre imágenes ópticas RGB o multiespectrales; para imágenes hiperespectrales, que registran cientos de bandas espectrales por píxel, su extensión constituye un problema poco explorado.

Llevar la detección semántica de cambios con vocabulario abierto al dominio hiperespectral abre un abanico de aplicaciones de alto impacto. Para esta aplicación en específico, se utilizan imágenes hiperespectrales bitemporales que corresponden a imágenes hiperespectrales de una misma escena en tiempos diferentes. La firma espectral de cada píxel actúa como una huella dactilar del material, de modo que un sistema capaz de detectar y describir cambios sobre este tipo de imágenes podría distinguir transformaciones imperceptibles para una cámara convencional. En el ámbito civil, esto resulta valioso para el monitoreo de la deforestación, la minería ilegal, la contaminación de cuerpos de agua, el crecimiento urbano, la salud de los cultivos y la respuesta ante desastres naturales, escenarios en los que la naturaleza del cambio importa tanto como su ubicación. En el ámbito militar y de seguridad, la riqueza espectral permite revelar materiales y actividades que el camuflaje visual oculta, como movimientos de vehículos, alteraciones del terreno o presencia de materiales artificiales, información difícil de obtener con sensores RGB tradicionales.

Para abordar esta brecha, el presente proyecto buscará integrar un conjunto de modelos fundacionales ya existentes en una única cadena de procesamiento: extracción de características espectrales y espaciales de las imágenes, localización de cambios, segmentación de regiones y asignación de etiqueta de vocabulario abierto, con la posibilidad de generar descripciones en lenguaje natural mediante modelos de generación de descripciones espectrales, o \textit{spectral captioning}.

Con base en lo anterior, este proyecto se orienta por la siguiente pregunta de investigación: ¿cómo realizar detección semántica de cambios con vocabulario abierto en imágenes hiperespectrales bitemporales mediante la alineación entre información espectral y textual? Responder a esta pregunta implica diseñar una solución que, sin entrenamiento supervisado específico para cada tipo de cambio, sea capaz de localizar, segmentar y describir con vocabulario abierto las diferencias entre dos adquisiciones hiperespectrales de una misma escena en tiempos diferentes.

El sistema propuesto se estructura como una cadena de procesamiento secuencial. Inicialmente se cargará el par de imágenes hiperespectrales de la misma escena en tiempos diferentes y se alinearán a un espacio espectral común \parencite{wang2024hypersigma}. Se extraerán representaciones profundas de cada imagen. Luego se calculará un mapa de cambio comparando ambas representaciones y se binarizará para aislar las regiones modificadas. Posteriormente, esas regiones se convertirán en indicaciones espaciales para producir máscaras precisas a nivel de píxel. Finalmente, se etiquetará cada máscara contra un vocabulario textual libre y, de forma opcional, se generará una descripción en lenguaje natural de cada región con un modelo de generación de descripciones espectrales \parencite{li2022blip,das2025hypercap}. Como resultado, el sistema buscará generar para cada par de imágenes un conjunto de regiones de cambio, cada una acompañada de su máscara, su etiqueta semántica y, cuando corresponda, una descripción textual del cambio detectado.

\pendiente{Alcance del trabajo: qué cubre y qué no cubre el sistema.}

\pendiente{Párrafo de cierre que anuncie la organización del documento, capítulo por capítulo.}
